% !TEX root = main.tex
% Section 2 — Related Work
% Five SOTA themes + explicit bridge paragraph.

\section{Related Work}
\label{sec:relatedwork}

This section synthesizes the predictive-BCI literature along five themes that emerged from a systematic review of 162 papers (2019--2026 with seminal anchors): deep sequence models for seizure prediction (\S\ref{sec:rw-sequence}), representation learning for neural signals (\S\ref{sec:rw-repr}), cross-patient generalization (\S\ref{sec:rw-xpatient}), efficient and deployable BCIs (\S\ref{sec:rw-efficient}), and cross-domain predictive monitoring (\S\ref{sec:rw-crossdomain}). A final bridge (\S\ref{sec:rw-bridge}) contrasts this body of work with the proposed study's commitments.

\subsection{Deep Sequence Models for Seizure Prediction}
\label{sec:rw-sequence}

Seizure prediction is the most methodologically mature subfield of predictive BCIs. Reviews \cite{WuEtAl2025,TanEtAl2025,RenEtAl2022,YimanWenqi2025,WongEtAl2023} consistently identify convolutional--recurrent hybrids and attention-based sequence models as the dominant architectural class, with CHB-MIT functioning as the field's near-universal benchmark.

Within this class, three recent results define the performance envelope. \citet{CaoEtAl2026} report 99.54\% AUC and a false-prediction rate of approximately 0.01/h on CHB-MIT using a multi-teacher knowledge-distillation model with 0.082M parameters --- an order of magnitude smaller than comparable architectures and explicitly aimed at deployment. \citet{AlkhrijahEtAl2025} achieve 99.34\% sensitivity via a feature-fusion ensemble that combines spectral, statistical, and deep-learned embeddings. \citet{PhamTran2026} extend dual-branch attention with a neural memory module, reporting 99.46\% sensitivity across long horizons and explicitly targeting horizon-stable prediction. Complementary work establishes the breadth of the architectural space: transformer variants \cite{MahdiBaghdadi2026,GodoyEtAl2022,ZhengEtAl2025,ShiLiu2023,WangEtAl2026a}, convolutional--recurrent hybrids \cite{DaoudBayoumi2019,ChoiEtAl2022,AbdelhameedBayoumi2021,HusseinEtAl2022,Chen2023,AnsariEtAl2020}, graph neural networks for spatial structure \cite{ZhouEtAl2026a,AmiriEtAl2026,AbbasiEtAl2025,LiaoEtAl2025a}, and causal or dynamical-systems-inspired designs \cite{QuEtAl2026,WangEtAl2025c,WangEtAl2025b,WangEtAl2026b} all report comparable within-dataset ceilings.

A second tier of work addresses specific technical bottlenecks: data scarcity via synthetic preictal generation \cite{XuEtAl2022,JiangEtAl2024a,CinquettiEtAl2025}, class imbalance via contrastive pretraining \cite{GuoEtAl2023,QiEtAl2025}, channel burden via channel-selection and single-channel models \cite{WangEtAl2021a,JangEtAl2026,WangEtAl2026a,Cherouati2023}, and interpretability via explainable wavelet and fuzzy-clustering designs \cite{WangEtAl2025b,WangEtAl2026b,SuffianEtAl2025}. Optimization-based variants --- firefly-trained deep networks \cite{PN2024}, Cramer-distance losses \cite{GhanimiEtAl2024}, heuristic-optimized meta-networks \cite{SEtAl2025}, memristive CNNs \cite{LiEtAl2022} --- round out a field whose architectural diversity is, if anything, saturating. A growing consensus \cite{WuEtAl2025,TanEtAl2025,RenEtAl2022,WongEtAl2023} is that within-dataset CHB-MIT performance is no longer the primary bottleneck.

\subsection{Representation Learning for Neural Signals}
\label{sec:rw-repr}

A parallel line of work treats the predictive task as secondary to the quality of the underlying neural representation. Self-supervised pretraining \cite{BanvilleEtAl2020,JiangEtAl2024b,BaryMacq2024} learns structure from unlabelled EEG before fine-tuning for downstream prediction. Geometry-aware encoders using symmetric positive-definite manifolds \cite{PaillardEtAl2025,ChenEtAl2026a} exploit the Riemannian structure of EEG covariance matrices. Contrastive objectives --- both within-subject \cite{GuoEtAl2023,QiEtAl2025} and cross-dataset \cite{LiaoEtAl2024} --- push representations toward invariance under nuisance variation. Generative models --- GANs \cite{XuEtAl2022,CinquettiEtAl2025,YinEtAl2026}, diffusion models \cite{JiangEtAl2024a,PuahEtAl2025}, and vision-transformer backbones fine-tuned from ImageNet \cite{YangModesitt2023,HoffsommerEtAl2025} --- serve both as augmentation engines and as probes of what structure the representation encodes. A smaller but conceptually important strand studies the \emph{predictive} content of neural representations themselves \cite{GunasekaranEtAl2025,ChenEtAl2026b,KölblEtAl2025,AkamaEtAl2025,BalarKapila2025,DercksenEtAl2026,EguinoaEtAl2026}, grounding the forecasting paradigm in the observation that the brain itself maintains temporal priors that can be read out.

Across this literature, the reported benefit of pretraining and representation learning is highest when downstream labelled data are scarce or when evaluation is cross-subject --- precisely the regime in which commodity supervised models degrade.

\subsection{Cross-Patient Generalization}
\label{sec:rw-xpatient}

The most diagnostic axis in the seizure-prediction literature is the gap between subject-specific and subject-independent performance. \citet{KoutsouvelisEtAl2024} quantify this gap directly: optimizing the preictal-period length and model choice yields a 20-percentage-point difference between the subject-specific regime (where each patient contributes training data) and the subject-independent regime (where models generalize to held-out patients). \citet{AmiriEtAl2026}, working with a physics-informed graph neural network grounded in chimera-state detection, report 89.3\% cross-patient sensitivity that decays to 68.2\% at a 90-minute forecasting horizon --- establishing a quantitative anchor for horizon-versus-generalization trade-offs.

Approaches to narrowing this gap fall into three families. \emph{Patient-independent training} \cite{DissanayakeEtAl2020,JangEtAl2026} pools data across patients and relies on architectural inductive biases or domain-adversarial objectives. \emph{Single-seizure or few-shot adaptation} \cite{TariqEtAl2020,JadduEtAl2024,MahdiBaghdadi2026} accepts that patient-level variation is irreducible and provides lightweight adaptation mechanisms at deployment. \emph{Cross-dataset validation studies} \cite{LiEtAl2025b,JangEtAl2026,MohammadSaeed2025,EsmaeilpourEtAl2024} characterize how specific design choices degrade across CHB-MIT, SNUH, MSSM, Siena, and related cohorts, typically reporting 3--15 percentage-point drops. Reviews \cite{WuEtAl2025,TanEtAl2025,RenEtAl2022} converge on the view that cross-patient generalization --- not within-dataset accuracy --- is the bottleneck for translation.

\subsection{Efficient and Deployable BCIs}
\label{sec:rw-efficient}

Deployment on wearables and smartphones imposes constraints orthogonal to accuracy: parameter budgets, memory footprint, latency, and energy. \citet{CaoEtAl2026}'s 0.082M-parameter architecture and \citet{AlkhrijahEtAl2025}'s ensemble stand at the efficient end of the Tier 1 spectrum. \citet{WangEtAl2021b} provide an enhanced benchmarking methodology for power-efficient prediction, and \citet{YedurkarEtAl2025}'s SPA-IoT + MCSV-CNN stack pairs efficient inference with an IoT alert channel. Memory-efficient adaptation techniques \cite{LiEtAl2025a,XuYu2026} explicitly target on-device fine-tuning, and lightweight statistical frameworks \cite{YangEtAl2025,WangEtAl2025b} eschew heavy deep-learning stacks for explainable, low-resource alternatives.

A second cluster of work couples efficient inference to IoT and smartphone infrastructure \cite{Paneru2026,KumarEtAl2026,ThahniyathEtAl2024,AbbasiEtAl2025}, treating the alert channel --- not only the predictor --- as a design concern. Edge-enabled geometric deep-learning approaches \cite{AbbasiEtAl2025} and cloud-based real-time analyses \cite{ThahniyathEtAl2024} illustrate the design space between full on-device and full cloud computation. Economic analyses of wearable deployment \cite{VelasquezEtAl2024} are beginning to pair cost-effectiveness questions with the technical feasibility demonstrated by this cluster. A general conclusion: running a Tier 1 seizure-prediction model on a wearable-paired smartphone is no longer the bottleneck to closed-loop intervention.

\subsection{Cross-Domain Predictive Monitoring}
\label{sec:rw-crossdomain}

Outside seizure epilepsy, the predictive paradigm recurs in adjacent adverse events, almost always independently of the BCI literature. \emph{Sleep} work forecasts non-apneic arousal from multichannel polysomnography with recurrent deep networks \cite{JamaraniEtAl2026} and, in adjacent work, forecasts survival or outcome from post-arrest EEG \cite{PelentritouEtAl2025,NesaragiEtAl2025,StaffordEtAl2026,ZhouEtAl2026b}. \emph{Migraine} prediction increasingly relies on wearable-derived autonomic and sleep features \cite{TomalinEtAl2025,KapustynskaEtAl2024,KapustynskaEtAl2025,JankevičiūtėEtAl2026}, re-inventing the preictal concept with different terminology. \emph{Fall} prediction in older adults couples wearable inertial sensing, gait biomarkers, and deep learning \cite{MaioraEtAl2024,ChanEtAl2023,Al-HammouriEtAl2025,CavanaghEtAl2024,GuanEtAl2025,LinEtAl2025,ZhangEtAl2024,LiaoEtAl2025b,ZuoEtAl2024,CesariEtAl2023}, and is increasingly the subject of scoping reviews and meta-analyses \cite{HrubýEtAl2026,MouEtAl2026,MehrlatifanMolla2024}. \emph{Vigilance and drowsiness} prediction in safety-critical settings \cite{TheivadasPonnan2025,SchwarzEtAl2023,Al-QuraishiEtAl2024,ZhangEtAl2025} shares the same computational skeleton --- preictal-like transition prediction --- under a different name. Cross-cutting reviews \cite{WesalEtAl2026,QamarEtAl2025,GongEtAl2025,Kumar2025} are beginning to articulate the multi-domain structure that this paper names explicitly.

The methodological asymmetry across domains is striking: seizure-epilepsy work contributes the bulk of Tier 1 evidence, while migraine, sleep, fall, and vigilance work --- despite its conceptual kinship --- lags in architectural sophistication, benchmark density, and the kind of saturated-within-dataset plateau that seizure prediction has reached. This asymmetry motivates our decision to locate the first preventive-outcome pilot in seizure epilepsy, where the predictive infrastructure is strongest, rather than in a domain where both the predictor and the intervention loop would be simultaneously immature.

\subsection{Bridge: From Prediction Quality to Preventive Outcome}
\label{sec:rw-bridge}

Taken together, the five themes above describe a field that has decisively solved the within-dataset prediction problem in seizure epilepsy, has made substantial progress on cross-patient generalization and on-device deployment, and is extending the predictive paradigm across adjacent adverse-event domains. What the field has not yet done --- and what this paper proposes to do --- is to \emph{measure} whether any of this prediction performance translates into a downstream preventive effect in an ambulatory human cohort. Almost every Tier 1 paper we surveyed reports prediction-quality metrics (sensitivity, false-prediction rate, AUC, lead time) on offline benchmarks; the reviews we cite \cite{WuEtAl2025,TanEtAl2025,RenEtAl2022,WesalEtAl2026,WongEtAl2023,YimanWenqi2025,CostaEtAl2024,QamarEtAl2025} note the absence of outcome-measurement studies explicitly, but no study in our corpus resolves it.

Our pilot is deliberately minimal along every axis except outcome measurement. The prediction model is a pre-validated CNN--LSTM--attention architecture of the class reported in \cite{CaoEtAl2026,AlkhrijahEtAl2025,LiEtAl2025b}, not a new contribution; the intervention protocol is standard clinical practice (behavioural safety checklist plus caregiver alert); the study population is restricted to drug-resistant focal epilepsy because it is the subfield with the densest Tier 1 prediction-quality evidence. The study's contribution is the within-subject pre/post measurement of seizure-rate reduction and the explicit linkage of that reduction to prediction lead time, sensitivity, and alert adherence --- the quantities that are pervasively reported as intermediate metrics but never, to our knowledge, linked to a downstream outcome. By isolating loop-closure as the causal contribution and deferring cross-domain extension, factorial generalization, and sham-controlled comparison to explicitly pre-specified follow-on studies, the pilot produces one unambiguous preventive-outcome estimate rather than a diffuse set of weak ones --- which we argue is the missing calibration point that the predictive-BCI field now most needs.
